% Deutsche Parallelfassung zu foundations/pose3_on_ellipsoid.tex

\chapter{Räumliche Pose in der Erdrealisierung}

Das vorherige Kapitel erzeugte die realisierte Flächentrajektorie,
\(\mathbf t\), \(\mathbf n_{\mathcal E}\), \(\mathbf b\) und
\(\kappa_g\). Dies ist keine zweite Alignment-Definition. Es sind die
metrischen Ausgaben, die zur Formulierung einer räumlichen Ingenieurfrage
benötigt werden: Wie wirken Überhöhung, Gravitation und Geschwindigkeit an
jeder Position desselben Alignments?

\section{Von der Flächentrajektorie zu pose3}

Die Anwendung des Frame-Operators ergibt das unüberhöhte Flächen-Frame
\[
F_{\mathcal E}(s)=
\bigl(\mathbf t(s),\mathbf b(s),\mathbf n_{\mathcal E}(s)\bigr).
\]
Die flächenbezogene Pose ist folglich
\[
\mathrm{pose3}_{\mathcal E}(s)=
\bigl(
\gamma_{\mathcal E}(s),
\mathbf t(s),
\mathbf n_{\mathcal E}(s),
\kappa_g(s),
\phi(s)
\bigr).
\]
Der ebene Richtungswinkel \(\theta\) wurde nicht willkürlich verworfen: Seine
Rolle wird nun von einem Tangentenvektorfeld übernommen, dessen Vergleich von
Punkt zu Punkt von der Flächenverbindung abhängt.

Die Krümmungsänderung ändert \(\kappa_g\) in dieser Realisierung unmittelbar
und damit die fortgepflanzte Tangente, die Flächentrajektorie und das lokale
Frame. Referenzellipsoid und Gravitationsmodell bleiben feste
Qualifikationen; die Überhöhung bleibt fest, sofern sie nicht bearbeitet oder
durch eine anwendbare Nebenbedingung gekoppelt wird.

\section{Überhöhung dreht den realisierten Querschnitt}

Die Überhöhungsdrehung \(\phi(s)\) wirkt um die Tangente:
\[
\mathbf b_{\phi}(s)
=R_{\mathbf t}(\phi(s))\,\mathbf b(s),
\qquad
\mathbf n_{\phi}(s)
=R_{\mathbf t}(\phi(s))\,\mathbf n_{\mathcal E}(s).
\]
Diese Operation empfängt das realisierte Flächen-Frame und den
Überhöhungszustand und erzeugt den überhöhten Querschnitt. Sie bewahrt die
Alignment-Identität und hält sowohl Zustands- als auch
Realisierungsprovenienz fest. Sie leitet weder Überhöhung aus einem gerenderten
Schienenmodell ab, noch entscheidet sie, dass das resultierende Gleichgewicht
akzeptabel ist.

\section{Gravitation macht die Pose physikalisch interpretierbar}

Es sei \(\mathbf g(s)\) der anwendbare lokale Gravitationsvektor an
\(\gamma_{\mathcal E}(s)\). Eine vereinfachte Realisierung kann
\(\mathbf g\parallel-\mathbf n_{\mathcal E}\) setzen; eine Realisierung
höherer Genauigkeit kann Ellipsoidnormale, Geoidnormale und lokale Gravitation
unterscheiden. Die Wahl gehört zum metrischen und physikalischen
Bewertungskontext, nicht zur Alignment-Identität.

Gravitation wird hier eingeführt, weil Überhöhung nur relativ zu einer
Kraftrichtung eine ingenieurtechnische Folge besitzt. Die Flächennormale
konstruiert das Frame; der Gravitationsvektor liefert die in der Bewertung
verwendete physikalische Richtung.

\section{Krümmung, Geschwindigkeit, Überhöhung und effektive Beschleunigung}

Für die Geschwindigkeit \(v\) lautet der krümmungsinduzierte Queranteil
\[
a_{\kappa}(s)=v^2\kappa_g(s).
\]
Die Projektion der Krümmungs- und Gravitationsanteile in die überhöhte
Seitenrichtung ergibt unter der verwendeten Vorzeichenkonvention
\[
a_{\mathrm{eff}}(s)
=v^2\kappa_g(s)
-\mathbf g(s)\cdot\mathbf b_{\phi}(s).
\]
Unter der lokalen ebenen Näherung reduziert sich dies auf
\[
a_{\mathrm{eff}}(s)
=v^2\kappa(s)-g\sin\phi(s).
\]

Die Abhängigkeit ist nun ausdrücklich: Die Krümmungsänderung ändert
\(\kappa_g\) und das realisierte Frame; die Geschwindigkeit liefert den
betrieblichen Zustand; die Überhöhung dreht den Querschnitt; die Gravitation
liefert die physikalische Richtung; die Bewertung erzeugt die effektive
Beschleunigung. Das Ergebnis hängt von allen fünf Eingaben ab. Es kann weder
eine zulässige Geschwindigkeit bestimmen, die Änderung freigeben noch die
Ingenieurregel ersetzen, unter der es beurteilt wird.

\section{Was diese Spezialisierung festlegt}

\begin{description}
\item[Heute nutzbar:] AIM kann Identität und Provenienz bewahren, während es
  einen intrinsischen Zustand rekonstruiert, einen ausdrücklichen metrischen
  Kontext anwendet, ebene Repräsentationen erzeugt und qualifizierte Folgen
  bewertet. Eine lokale ebene Realisierung bleibt ausreichend, wann immer
  ihre deklarierte Genauigkeit für die Aufgabe genügt.
\item[Formale Spezialisierung:] Diese Kapitel formulieren, wie sich dieselben
  Operatorverträge auf ein Referenzellipsoid spezialisieren, einschließlich
  geodätischer/normaler Zerlegung, Flächenpose, überhöhtem Frame und
  gravitationsbewusster Abhängigkeit. Dies ist eine mathematische
  Spezialisierung, keine Behauptung fertiggestellter Produktionssoftware.
\item[Forschung:] Validierte Anwendbarkeitskriterien, Fehlergrenzen zwischen
  ebenen/projizierten und ellipsoidischen Realisierungen, robuste numerische
  Realisierung über Daten im Projektmaßstab und die erforderliche Genauigkeit
  von Gravitationsmodellen bleiben Forschungs- und Validierungsarbeit.
\item[Praktischer Umfang:] AIM verlangt weder, jede Änderung auf einem
  Ellipsoid zu berechnen, noch etablierte Projekt-CRS-Arbeitsabläufe zu
  ersetzen oder einen geodätischen Solver einzusetzen, wenn eine qualifizierte
  ebene Realisierung die Ingenieurfrage beantwortet.
\end{description}

Erdkrümmung hat damit einen begrenzten und genauen Platz in der formalen
Maschinerie erhalten. Sie wird notwendig, wenn der Realization Context
Flächeneffekte wesentlich macht; andernfalls bleibt sie eine verfügbare
Spezialisierung. In beiden Fällen bleiben Alignment-Identität, konstruktive
Krümmungsgeschichte und Verantwortungsgrenzen erhalten.
